\documentclass[11pt,a4paper]{article}
\usepackage[margin=1in]{geometry}
\usepackage[T1]{fontenc}
\usepackage[utf8]{inputenc}
\usepackage{lmodern}
\usepackage{enumitem}
\usepackage{hyperref}
\usepackage{xcolor}
\usepackage{array}
\usepackage{longtable}
\usepackage{setspace}

\setlength{\parskip}{0.5em}
\setlength{\parindent}{0pt}

\begin{document}

% ----------------------------
% Front Page
% ----------------------------
\begin{titlepage}
	\centering
	\vspace*{2.0cm}

	{\Large \textbf{Indian Institute of Technology Delhi}}\\[0.6cm]
	{\large \textbf{COP290: Design Practices in Computer Science}}\\[1.0cm]

	{\Huge \textbf{Assignment 4 Report}}\\[0.35cm]
	{\LARGE \textbf{DataFrameLib Implementation}}\\[2.0cm]

	\renewcommand{\arraystretch}{1.5}
	\begin{tabular}{>{\bfseries}p{5.0cm} p{8.0cm}}
		Student Name & Megh Bhavesh Jesalpura \\
		Entry Number & 2024CS10844 \\
		Course Code & COP290 \\
		Assignment & A4 \\
		Submission Date & \today \\
	\end{tabular}

	\vfill
	\textit{This report documents the implementation choices, algorithmic details, and correctness arguments for the EagerDataFrame, LazyDataFrame, and QueryOptimizer components based on the submitted codebase.}
\end{titlepage}

% ----------------------------
% Contents Page
% ----------------------------
\pagenumbering{roman}
\tableofcontents
\newpage

\pagenumbering{arabic}
\onehalfspacing

\section{EagerDataFrame and Exactly How Different Strategies Have Been Implemented}

\subsection{Overview of Eager Execution Philosophy}
The eager engine is built around a direct and deterministic model: each API call executes immediately on an in-memory Arrow table and returns a fresh \texttt{EagerDataFrame}. This has two practical consequences in the implementation. First, there is no deferred state to maintain; each object is always materialized and fully usable. Second, operation-level errors (missing columns, invalid predicate types, unsupported join kinds) are surfaced at the same call site instead of at a later \texttt{collect()} stage.

The class stores only one private member, \texttt{std::shared\_ptr<arrow::Table> table\_}, and every method reconstructs an output table explicitly from Arrow arrays and schema fields. This keeps ownership simple and makes memory behavior easier to reason about.

\subsection{Input/Output and Construction Strategy}

\textbf{CSV and Parquet loading.} \texttt{read\_csv()} opens a readable file handle and creates Arrow's CSV table reader with default read/parse/convert options. The resulting Arrow table is returned immediately after status checks. \texttt{read\_parquet()} opens a Parquet file using \texttt{parquet::arrow::OpenFile}, reads the entire table, and wraps it in \texttt{EagerDataFrame}.

\textbf{From-columns API strategy.} Two construction paths are provided:
\begin{itemize}[leftmargin=*]
	\item \texttt{from\_columns(map<string, ChunkedArray>)} when data already exists in Arrow chunked form.
	\item \texttt{from\_columns(vector<pair<string, Array>>)} for ordered insertion while wrapping each array into a one-chunk \texttt{ChunkedArray}.
\end{itemize}
This dual API is useful because internal operations often produce flat arrays while external callers may already use chunked data.

\textbf{Output strategy.} CSV writes are performed through \texttt{WriteCSV} after \texttt{CombineChunksToBatch()}, and Parquet writes use \texttt{parquet::arrow::WriteTable}. Both methods throw detailed runtime errors when status objects indicate failure.

\subsection{Strategy for Projection, Filtering, and Derived Columns}

\textbf{Projection (\texttt{select}).} The implementation loops through requested names, resolves both field and column by name, and builds a new schema plus array list. Missing names trigger exceptions immediately. Complexity is approximately linear in number of selected columns and independent of row count, since it mostly reuses existing arrays.

\textbf{Filtering (\texttt{filter}).} The method evaluates an expression to a mask and enforces that the expression's result type is boolean. A key strategy here is explicit mask application column-by-column:
\begin{enumerate}[leftmargin=*]
	\item Evaluate predicate to a boolean chunked array.
	\item Flatten mask for contiguous row access.
	\item For each column, flatten and apply \texttt{applyBooleanMask(...)}.
	\item Rebuild a new table with identical schema and filtered arrays.
\end{enumerate}
This avoids hidden behavior in third-party filtering APIs and gives full control over null handling and type dispatch in utility functions.

\textbf{Column creation/replacement (\texttt{with\_column}).} The expression is evaluated once, then existing column with same name (if present) is skipped during reconstruction. The new column and its field are appended. This realizes a ``replace-if-exists, append-if-new'' policy.

\subsection{Strategy for Head and Sorting}

\textbf{Head.} \texttt{head(n)} slices each column using Arrow \texttt{Slice(0, min(n, length))}. Since slicing is per-column, row count in all output columns remains consistent.

\textbf{Sorting strategy details.} Sorting has been implemented with a precomputation-first approach:
\begin{itemize}[leftmargin=*]
	\item Sort-key columns are flattened once and converted into in-memory typed vectors.
	\item Numeric values are promoted to \texttt{double} in key storage to unify comparison logic.
	\item Null masks are stored explicitly per sort key and compared before value comparisons.
	\item \texttt{std::stable\_sort} is used to preserve deterministic ordering for ties.
	\item Post-sort reordering is performed for all columns using a single sorted index list.
\end{itemize}
This strategy shifts work from the comparator into preprocessing, reducing overhead during \(O(N \log N)\) compare calls.

\subsection{Join Strategy in Depth}
The join implementation is a hash-join with practical engineering choices for performance and correctness:

\textbf{1) Build side selection.} The smaller relation is chosen as hash build side to reduce memory footprint and improve cache locality.

\textbf{2) Key encoding.} Composite keys are serialized to compact binary strings. For each key component:
\begin{itemize}[leftmargin=*]
	\item nulls are marked with a dedicated sentinel byte,
	\item non-null numeric values are appended as raw bytes,
	\item strings are appended directly,
	\item a delimiter byte separates components.
\end{itemize}
This avoids repeated scalar-to-string conversion for numeric keys.

\textbf{3) Probe and match collection.} Probe rows generate keys and lookup in hash map \texttt{key -> vector<row\_indices>}.

\textbf{4) Outer semantics.} For \texttt{left}, \texttt{right}, \texttt{outer} joins, unmatched tuples are explicitly appended with \texttt{-1} for missing side index.

\textbf{5) Output schema strategy.} Join keys are coalesced into one key column in output. Non-key right columns are appended, and conflicts are disambiguated by \texttt{\_right} suffix.

\subsection{GroupBy and Aggregation Strategy in Depth}

\textbf{Grouping state representation.} Grouping first pre-flattens all group-key arrays and numeric candidate aggregate arrays. Group identity is encoded as binary key and aggregated through \texttt{unordered\_map} for expected \(O(1)\) updates.

\textbf{Accumulator design.} For each (group, column), \texttt{GroupAccumulator} stores:
\[
\text{sum},\; \text{count},\; \text{min},\; \text{max}
\]
so all supported aggregate outputs can be derived without rescanning data.

\textbf{Finalization strategy.} Aggregation outputs are materialized column-wise, with output Arrow type selected per operation and source type (for example, \texttt{count} always as \texttt{int64}).

\subsection{Complexity Summary for Eager Operations}
\begin{longtable}{|p{3.1cm}|p{5.4cm}|p{5.0cm}|}
\hline
\textbf{Operation} & \textbf{Main Strategy} & \textbf{Asymptotic Cost (rough)} \\
\hline
\texttt{select} & Schema+array reconstruction from named columns & \(O(k)\) for \(k\) selected columns \\
\hline
\texttt{filter} & Evaluate mask then apply to each column & \(O(R \cdot C)\) \\
\hline
\texttt{with\_column} & Evaluate expression once and rebuild schema/arrays & \(O(R + C)\) \\
\hline
\texttt{sort} & Precompute keys + stable sort + reorder columns & \(O(R \log R + R\cdot C)\) \\
\hline
\texttt{join} & Hash join with smaller-side build & Expected \(O(R\_l + R\_r + M)\) \\
\hline
\texttt{group\_by+aggregate} & Hash grouping + accumulator finalization & Expected \(O(R\cdot G)\) \\
\hline
\end{longtable}

\section{LazyDataFrame and Exactly How I Am Making a Graph Using std::variant}

\subsection{Why std::variant Was Chosen}
The logical plan is represented as a value-type sum over node structs:
\[
\texttt{planNode = variant<ScanNode, FilterNode, SelectNode, WithColumnNode, GroupByNode, AggNode, JoinNode, SortNode, HeadNode>}
\]
This gives compile-time exhaustiveness with \texttt{std::visit} and removes the need for a base-class hierarchy with virtual dispatch for plan nodes.

\subsection{Plan Node Data Model}
Every node contains exactly the parameters needed for that relational operator:
\begin{itemize}[leftmargin=*]
	\item \texttt{ScanNode}: file path + CSV/Parquet flag.
	\item Unary operators: payload + \texttt{child} pointer.
	\item \texttt{JoinNode}: \texttt{left}, \texttt{right}, join keys, join type.
	\item \texttt{HeadNode}: row limit \(n\) + child.
\end{itemize}
The data model is explicit and serializable to DOT labels easily.

\subsection{How Lazy Operations Build the DAG}
Each chained API call wraps the old root in a new node:
\begin{itemize}[leftmargin=*]
	\item \texttt{scan\_csv(path)} starts with \texttt{ScanNode(path, is\_csv=true)}.
	\item \texttt{filter(pred)} creates \texttt{FilterNode(pred, oldRoot)}.
	\item \texttt{select(cols)} creates \texttt{SelectNode(cols, oldRoot)}.
	\item \texttt{with\_column(name, expr)} creates \texttt{WithColumnNode(name, expr, oldRoot)}.
	\item \texttt{group\_by(keys).aggregate(aggs)} forms \texttt{GroupByNode} then \texttt{AggNode}.
	\item \texttt{join(other,...)} creates binary \texttt{JoinNode(oldRoot, otherRoot,...)}.
\end{itemize}
No data is touched during these calls. Only node objects and pointers are constructed.

\subsection{From DAG to Execution: Variant-Visitor Runtime}
\texttt{collect()} follows this pipeline:
\begin{enumerate}[leftmargin=*]
	\item optimize logical plan,
	\item apply projection pushdown pass,
	\item execute optimized plan recursively.
\end{enumerate}

Execution uses \texttt{std::visit} where each node case maps directly to eager operator calls. This makes semantics transparent: plan nodes are merely delayed forms of eager operations.

\subsection{Exactly How Graph Rendering Is Implemented}
The graph rendering path is deterministic and two-step:

\textbf{Step 1: DOT generation.}
\begin{itemize}[leftmargin=*]
	\item \texttt{LazyDataFrame::explain(path)} calls optimizer first.
	\item It opens \texttt{path + ".dot"} and writes graph preamble (direction, node style).
	\item It calls recursive helper \texttt{writeDot(out, optimized, nodeId)}.
\end{itemize}

\textbf{Step 2: PNG rendering.}
\begin{itemize}[leftmargin=*]
	\item System command: \texttt{dot -Tpng <dotfile> -o <path>}.
	\item Non-zero return value raises runtime error.
\end{itemize}

\textbf{Inside \texttt{writeDot}:}
\begin{itemize}[leftmargin=*]
	\item A unique integer \texttt{myId} is assigned to each visited node.
	\item \texttt{std::visit} pattern-matches node type and prints a readable label:
	\begin{itemize}[leftmargin=*]
		\item Scan labels include file path,
		\item Filter labels include predicate expression string,
		\item Select/GroupBy/Agg labels include column/operator payload,
		\item Join labels include join type and two child edges,
		\item Sort labels include order direction,
		\item Head labels include row limit \(n\).
	\end{itemize}
	\item For child nodes, recursion returns child ids used to write \texttt{parent -> child} edges.
\end{itemize}

So the variant-based DAG is not only used for execution but also for user-facing plan explainability.

\subsection{Correctness Intuition for Lazy-to-Eager Mapping}
For each node type, execution rule in \texttt{QueryExecutor} is exactly the eager counterpart. Therefore, for a fixed plan \(P\), the result of lazy \texttt{collect()} equals eager composition of the same operations in the same tree shape (modulo safe optimizer rewrites).

\section{QueryOptimizer and Correctness of the Transformations That I Have Made}

\subsection{Optimizer Architecture}
The optimizer has two public static entry points:
\begin{itemize}[leftmargin=*]
	\item \texttt{optimize(planNode input)}: bottom-up rewrites and local transformation rules.
	\item \texttt{pushdownProjections(const planNode\& input)}: dependency-driven column pruning pass.
\end{itemize}

The first pass simplifies expressions and pushes operators where safe. The second pass computes required columns and rewrites plan to avoid unnecessary column flow through the DAG.

\subsection{Transformation A: Constant Folding and Algebraic Simplification}

\textbf{What is implemented.}
The optimizer detects constant-only expressions by checking absence of \texttt{col(...)} references in \texttt{toString()}. Such expressions are evaluated on a one-row dummy table and replaced with literal expressions. In addition, recursive simplifications are applied:
\begin{itemize}[leftmargin=*]
	\item \(x + 0 = x\), \(0 + x = x\), \(x - 0 = x\)
	\item \(x \cdot 1 = x\), \(1 \cdot x = x\)
	\item \(x \cdot 0 = 0\), \(0 \cdot x = 0\)
	\item \(x / 1 = x\)
	\item \(\text{true} \land x = x\), \(\text{false} \land x = \text{false}\)
	\item \(\text{false} \lor x = x\), \(\text{true} \lor x = \text{true}\)
	\item \(\neg(\neg x) = x\)
\end{itemize}

\textbf{Correctness argument.}
All rewrites are tautological identities in arithmetic/boolean algebra. Replacing subtree \(E\) by equivalent \(E'\) preserves row-wise evaluation for every input tuple. Since relational operators consume only these values, table output remains unchanged.

\subsection{Transformation B: Predicate Pushdown}

\textbf{Rules implemented.}
\begin{itemize}[leftmargin=*]
	\item Filter through Select: always safe.
	\item Filter through Sort: always safe.
	\item Filter through WithColumn: safe only if predicate does not reference newly created column token.
	\item Filter through Join: side-aware and join-type-aware conditions.
\end{itemize}

\textbf{Join-aware safety used in code.}
\begin{itemize}[leftmargin=*]
	\item \texttt{inner}: push to either side if predicate references only that side.
	\item \texttt{left}: push only to right side.
	\item \texttt{right}: push only to left side.
	\item \texttt{outer}: no push.
\end{itemize}

If predicate columns are all join keys, the code may push to both sides for applicable join kinds.

\textbf{Correctness argument.}
For unary row-preserving operators (Select, Sort), predicate commutation is straightforward. For WithColumn, the guard avoids changing dependencies. For joins, restrictions preserve null-extension behavior of outer joins and therefore avoid cardinality/row-content regressions.

\subsection{Transformation C: Limit Pushdown}
Implemented rules:
\[
\texttt{Head(Select(X)) -> Select(Head(X))},\quad
\texttt{Head(WithColumn(X)) -> WithColumn(Head(X))}
\]
No pushdown through operators where semantics may change (sort/group/aggregate/join/filter).

\textbf{Correctness argument.}
Both Select and WithColumn preserve row order and count. Therefore, taking first \(n\) rows before or after these operators yields identical row set and order.

\subsection{Transformation D: Projection Pushdown}
The second pass recursively propagates \texttt{requiredCols}:
\begin{itemize}[leftmargin=*]
	\item Filter extends requirements with predicate-referenced columns.
	\item Sort extends with sort keys.
	\item WithColumn is removed when produced column is never used.
	\item Agg trims aggregation list to needed outputs and pushes source requirements.
	\item Join splits requirements into left/right subsets plus mandatory join keys.
\end{itemize}

\textbf{Important practical note from current code.}
At \texttt{ScanNode}, no physical file-level column pruning is done yet because eager scan currently reads full tables. So this pass reduces intermediate column flow in memory but does not yet reduce scan I/O.

\textbf{Correctness argument.}
If a column is removed, it means no downstream operator needs it either directly or indirectly. Dependency closure is preserved by adding all operator-required columns (predicates, sort keys, aggregate inputs, join keys). Therefore, removed columns are observationally irrelevant.

\subsection{End-to-End Correctness Statement}
Let \(P\) be original plan and \(P'\) be optimized plan. Each rewrite in the implemented set is semantics-preserving under its explicit guard conditions. Since rewrites are composed sequentially and each step preserves semantics, final plan preserves semantics:
\[
\forall D,\; \texttt{execute}(P, D) = \texttt{execute}(P', D)
\]
where \(D\) is any valid input dataset satisfying schema/type assumptions.

\subsection{Performance Impact Summary}
Expected gains from implemented transforms:
\begin{itemize}[leftmargin=*]
	\item Less expression evaluation work via simplification/folding.
	\item Fewer rows entering expensive operators via predicate pushdown.
	\item Lower column movement overhead in intermediate nodes via projection pass.
	\item Earlier cardinality reduction via safe head pushdown.
\end{itemize}
Observed benefit magnitude depends on selectivity, width of schema, and operator order in user query chains.

\section{Design Decisions (Placeholder)}

\textit{Placeholder: Add final design rationale write-up here. Suggested subtopics to fill later:}
\begin{itemize}[leftmargin=*]
	\item \textit{Why Arrow-native table representation was preferred over custom column containers.}
	\item \textit{Trade-offs of variant-based plan nodes versus classical inheritance.}
	\item \textit{Memory ownership policy with \texttt{shared\_ptr} and how it affects copy behavior.}
	\item \textit{Rationale for hash-based join/grouping key encoding choices.}
	\item \textit{Future extension plan for physical projection pushdown at scan level.}
\end{itemize}

\section{Conclusions}

This project now has a coherent two-engine architecture: an eager execution core over Arrow tables and a lazy query-planning layer over a variant-based DAG. The eager side is implementation-driven and explicit, with stable ordering in sort, hash-based join, and accumulator-based grouping. The lazy side captures query intent and supports explainability through DOT/Graphviz rendering of optimized plans.

From an optimization perspective, the current rule set is already meaningful: expression simplification reduces scalar work, predicate/limit pushdown reduce row flow where safe, and projection propagation reduces unnecessary column movement across operators. The correctness guards in join-aware pushdown and with-column dependency checks are critical and are already present in the implementation.

The next natural milestone is to connect logical projection pushdown to physical scan pruning so that CSV/Parquet readers load only required columns. Once that is added, the optimizer can improve both CPU and I/O costs, making the lazy path significantly more scalable on wide datasets.

Overall, the submitted implementation satisfies the intended A4 direction: clear eager semantics, composable lazy DAG construction with \texttt{std::variant}, and optimizer transformations with explicit correctness intent.

\end{document}
